
\chapter{Screen scraping}

Welcome to the worst of all possible worlds\footnote{Actually, the
\textit{second}-worst; the very worst is when the data you seek isn't even in
electronic form at all.}: when the data you want to analyze is on a
human-readable but machine-unfriendly website somewhere. It's there, you can
see it...but getting it from a pretty web page into a Pandas DataFrame is a
non-trivial proposition.

Usually this situation occurs when the provider of the data, and the designer
of the website, didn't envision anyone using the data for any purpose other
than viewing it through the website itself. Sometimes there is another reason:
that the organization wanted to deliberately make this operation difficult so
that you \textit{couldn't} easily use it for other purposes. Either way,
though, as long as the data is accessible and visible through your browser, it
is at least theoretically extractable into the form we need to do our magic.

\section{HTML}

(Most) web pages look snazzy and colorful when you view them: the word
``code''
doesn't jump immediately to mind. When you load a web page, though, what
you're actually getting from the server you specify is a plain-text file
written in a language called HTML, or \textbf{hypertext markup language}. It
consists of a combination of content and formatting/structural commands (or
``markup'') to instruct the browser how to \textbf{render} the page.
(``Rendering'' a page means ``interpreting all that HTML formatting stuff and
actually presenting a pretty page formatted the way it specifies.)

The first important thing to know about HTML is that it is
\textit{hierarchical}, not flat. This makes it more like JSON than CSV. Even
though every HTML file is a single sequence of characters, it encodes a
tree-like structure in which elements can exist inside other elements, which
can exist inside other elements, \textit{etc.}, with no limit to the amount of
nesting.

The basic building block of HTML is the \textbf{element}, which is delimited
by starting and ending \textbf{tags}. Tags are key words, most of which
indicate something structural or stylistic, enclosed in wakkas (``\texttt{<
... >}''). The ending tag has a slash (``\texttt{/}'') immediately before the
keyword, showing that it matches the corresponding opening tag. So, for
instance, one pair of starting and ending tags might be ``\texttt{<title>}''
and ``\texttt{</title>}''.

Note that the entire contents of the \textit{element} is (1) the starting tag,
(2) the ending tag, and (3) \textit{everything in between}, which could be
text in addition to other elements.

\subsection{An example HTML file}

Figure~\ref{fig:allAboutStephen} shows a very simple, but complete, example.
Lots going on there. The first way to get your bearings is to start visually
lining up starting tags with their corresponding ending tags and perceive the
structure. Notice that this entire page begins with ``\texttt{<html>}'' and
ends with ``\texttt{</html>}''. This is in essence the ``root node'' of the
page's entire hierarchical tree. It has two immediate ``children'': a
\texttt{head} element, and a \texttt{body} element. Each of those have
starting and ending tags as well. Inside the \texttt{head} element is a
\texttt{title} element, \textit{etc.}

\begin{figure}
\centering
\begin{Verbatim}[fontsize=\scriptsize,samepage=true,frame=lines]
<html>
<head>
    <title>All About Stephen</title>
    <style>
        img { width:150px; }
        td { border-width:1px; }
    </style>
</head>
<body>
<h1>Stephen Davies, Esquire</h1>

<p>Deep in the bowels of
<a href="http://buildings.umwblogs.org/trinkle-hall/">Trinkle Hall</a> resides a
carbon-based life form known to outsiders only as <b>Stephen</b>. This
mysterious being fulfills several roles <u>not</u> widely known to
humankind.</p>

<p><i>Some</i> of these functions are delineated below. Please memorize them
and then immediately destroy this page upon reading it.</p>

<h2>Areas of operation</h2>

<table border=2 cellpadding=15>
    <tr><td class="pic"><img src="http://clipart.com/prof.jpg" width=200 /></td>
        <td class="role">College professor</td>
        <td class="description">Stephen teaches computer science and data science
    to willing undergraduates, shaping them into the heroes of tomorrow. </td>
        <td><ul><li>Super rewarding</li>
            <li>Exciting</li>
            <li>Ever-changing</li>
            <li>Fun</li></ul></td>
        <td><b>50%</b> of total time</td>
    </tr>
    <tr><td class="pic"><img src="http://clipart.com/dad.png" width=200 /></td>
        <td class="role">Dad</td>
        <td class="description"> With a wife, three teenagers, and two cats (one
    of whom throws up a lot), Stephen's home life is a constant source of love,
    anxiety, and Hallmark moments.</td>
        <td><ul><li>Amazing</li>
            <li>Loving</li>
            <li>Complicated</li>
            <li>Surprising</li></ul></td>
        <td><b>30%</b> of total time</td>
    </tr>
</table>
</body>
</html>
\end{Verbatim}
\caption{A sample HTML page.}
\label{fig:allAboutStephen}
\end{figure}

As you can guess, each of these tags is interpreted by the browser to mean a
particular kind of structure/formatting. ``\texttt{h1}'' means ``top-level
heading'' which in practicality means ``render it in a large font, with some
extra space around it.'' The ``\texttt{b}'' element means to put something in
boldface. The ``\texttt{table}'' element means to construct a table-like
structure, with rows (``\texttt{tr}'') and cells within rows (``\texttt{td}'').
\textit{Etc.} HTML is in general a big mess of things that are supposed to
indicate structure (like \texttt{table} and \texttt{h2}) and things that
specify specific formatting (like \texttt{i} and \texttt{u}). All of this is
eminently look-up-able online.

\pagebreak
Some starting tags (not ending) contain additional information called
\textbf{attributes}. You can see this with the \texttt{a} and \texttt{img}
tags from the example:

\begin{Verbatim}[fontsize=\footnotesize,samepage=true,frame=lines]
<a href="http://buildings.umwblogs.org/trinkle-hall/">Trinkle Hall</a>
\end{Verbatim}

The \texttt{a} tag stands for ``anchor'' and is the most common way of
embedding a hyperlink (which you can click on to go to another page) in a
document. Note that the words that will appear in the page are the
\textbf{contents} of the element, not part of the tag itself. The tag, in
addition to the required ``\texttt{a}'' has a key/value pair where the key is
the string \texttt{href} and the value is the URL the user will be directed to
if they click on the link. This key/value pair is called an
\textbf{attribute}.

\begin{Verbatim}[fontsize=\footnotesize,samepage=true,frame=lines]
<img src="http://clipart.com/prof.jpg" width=200 />
\end{Verbatim}

In the \texttt{img} example, we have \textit{two} attributes: one for the
filename of the image, and the other for some specific formatting suggestions
(in this case, a width of 200 pixels). An element can have many different
attributes, in any order.

(By the way, the sharp-eyed reader may have noticed something else odd about
the \texttt{img} example, above: an extra slash before the ending
``\texttt{>}''. This is an HTML shorthand for ``starting and ending the element
all within a single tag, since it has no content anyway.'' Whenever you have
something of the form ``\texttt{<x></x>}'' you can replace it with the
equivalent ``\texttt{<x/>}'' to save typing and file size.)

Lastly, I'll point out that near the top of this example, there's a bunch of
stuff which looks a bit visually jarring; it doesn't look like HTML. And it's
not. I'm referring to this:

\begin{Verbatim}[fontsize=\footnotesize,samepage=true,frame=lines]
    <style>
        img { width:150px; }
        td { border-width:1px; }
    </style>
\end{Verbatim}

This code is called \textbf{CSS} (or \textbf{Cascading Style Sheets}) which is
sort of a companion technology to HTML to give additional formatting control.
(Originally, HTML was supposed to strictly delineate document
\textit{structure}, and CSS could then be used to specify the formatting of
that structure. This strict separation of concerns, which was a good idea, has
gotten very muddy in the intervening years, however.)

\pagebreak
The whole thing, once rendered, looks something like this:

\begin{center}
\includegraphics[width=0.9\textwidth]{stephenpage.pdf}
\end{center}

\section{Web Developer Tools}

Firefox, Chrome, and several other browsers come equipped with very
sophisticated ``web developer tools'' that let you browse the HTML, CSS, and
rendered page, and figure out what elements got rendered via which portions of
the screen. They're very handy, and you should get into the habit of using
them. In Firefox, they're available via the ``Tools | Web Developer | Toggle
tools...'' menu option. In Chrome, you go to the little three vertical dots,
then choose ``More Tools | Developer Tools.''


\section{Parsing with BeautifulSoup}

From the point of view of Data Science, we don't care very much about all this
formatting stuff. To us, it's clutter. On the page is some important data we'd
like to analyze, and the goal is to sift through all this HTML and CSS crapola
looking for it and extracting it. To help us with this task is a Python
library called BeautifulSoup.

I generally recommend first \textit{downloading} the web page you want to
parse rather than having Python read it for you over the network. This
eliminates one more possible source of error. (Exception: if you need a Python
program to act as a \textbf{web crawler} and read pages on demand as it
discovers their URLs. More on that later.) Note that you can still open the
page in your browser even after downloading it: just go to ``File | Open
File...'' or the equivalent operation in your browser.

Once you've done this, and now have your HTML in a file called
``\texttt{myfile.html}'', then you can parse it with BeautifulSoup using:

\begin{Verbatim}[fontsize=\small,samepage=true,frame=single]
from bs4 import BeautifulSoup
soup = BeautifulSoup(open("myfile.html"), "html.parser")
\end{Verbatim}

By convention, the variable for this parsed structure is called
``\texttt{soup}''. The second argument to the \texttt{BeautifulSoup} function
tells it which \textbf{parser} we want to use to analyze the code; there are
other choices here, but stick with \texttt{"html.parser"} for now.

One possibly useful command (though limited if your web page is large) is:

\begin{Verbatim}[fontsize=\small,samepage=true,frame=single]
>>> print(soup.prettify())
\end{Verbatim}

This prints out the contents of the page, but in a way that is ``properly''
structured and indented. You may have noticed that the HTML file may or may
not be physically indented and spaced in a way reflective of the structure;
the parser discerns the implied structure from the tag sequence, and prints it
how it will ``really'' be seen by a browser.

At this point, I generally take one of two approaches: either I navigate the
structure \textbf{top-down} looking for where the data I'm interested in is
situated, or I find the data by searching through the file, and use a
\textbf{bottom-up} approach to get it from there.

\subsection{The top-down approach}

By using ``\texttt{soup.contents}'', you get the immediate descendants of
the root node of the HTML hierarchy. Each of these could be examined for its
type, like this:

\vspace{-.1in}
\begin{Verbatim}[fontsize=\footnotesize,samepage=true,frame=single]
>>> for child in soup.contents:
        print(type(child))

    <class 'bs4.element.Tag'>
    <class 'bs4.element.NavigableString'>
\end{Verbatim}

So our entire document is composed of an element (misleadingly named ``Tag''
here, but w/e) and a ``\texttt{NavigableString}'' which essentially just means
``string of literal text''. For elements, it's useful to print out their
``\texttt{.name}''. For strings, it's useful to print them out verbatim. Here,
we see:

\vspace{-.1in}
\begin{Verbatim}[fontsize=\footnotesize,samepage=true,frame=single]
>>> soup.contents[0].name
    'html'
>>> soup.contents[1]
    '\n'
\end{Verbatim}

So the first child is the entire \texttt{html} element in all its embedded
glory, and this is followed by a single, pesky newline. Now you know.

Drilling down another level, we see:

\vspace{-.1in}
\begin{Verbatim}[fontsize=\footnotesize,samepage=true,frame=single]
>>> html_children = soup.contents[0]
>>> for child in html_children:
        print(type(child))
    <class 'bs4.element.NavigableString'>
    <class 'bs4.element.Tag'>
    <class 'bs4.element.NavigableString'>
    <class 'bs4.element.Tag'>
    <class 'bs4.element.NavigableString'>
\end{Verbatim}

Two tags, with three strings interleaved. Let's examine them:

\vspace{-.1in}
\begin{Verbatim}[fontsize=\footnotesize,samepage=true,frame=single]
>>> html_children.contents[0]
    '\n'
>>> html_children.contents[1].name
    'head'
>>> html_children.contents[2]
    '\n'
>>> html_children.contents[3].name
    'body'
>>> html_children.contents[4]
    '\n'
\end{Verbatim}

You're starting to see a pattern: a lot of the strings are just newline
characters that split up the HTML text file and make it more readable. What we
really care about is the content under the tags. Let's dig one deeper:

\vspace{-.1in}
\begin{Verbatim}[fontsize=\footnotesize,samepage=true,frame=single]
>>> head_children = html_children.contents[1].contents
>>> head_children[0]
    '\n'
>>> head_children[1].name
    'title'
>>> head_children[2]
    '\n'
>>> head_children[3].name
    'style'
>>> head_children[4]
    '\n'
\end{Verbatim}

    Okay, you get the idea. We can drill down to whatever level we want just
by following the appropriate children. To get the title of the page, we could
do:

\begin{Verbatim}[fontsize=\footnotesize,samepage=true,frame=single]
>>> head_children[1].contents[0]
    'All About Stephen'
\end{Verbatim}

(This is because the second element (element \#1) of the \texttt{head} element
is the \texttt{title} element, and its only child is the string \texttt{'All
About Stephen'}.)

Sometimes a quicker method to find your way through the forest is to use this
syntax at the interactive console:

\begin{Verbatim}[fontsize=\footnotesize,samepage=true,frame=single]
>>> [ x.name for x in body_children ]
    [None, 'h1', None, 'p', None, 'p', None, 'h2', None, 'table', None]
\end{Verbatim}

You can better see at a glance what's what. Now that I know the first
``\texttt{p}'' is element \#3 of the body, I could find that hyperlink via:

\begin{Verbatim}[fontsize=\footnotesize,samepage=true,frame=single]
>>> [ x.name for x in list(body_children)[3] ]
    [ None, 'a', None, 'b', None, 'u', None]
\end{Verbatim}

and grab it via:

\begin{Verbatim}[fontsize=\footnotesize,samepage=true,frame=single]
dat_hyperlink =  list(list(body_children)[3])[1]
\end{Verbatim}

BeautifulSoup lets you access the attributes of an HTML element by treating
the element as a dictionary. Pretty cool. I can snag the actual link, then,
with:

\begin{Verbatim}[fontsize=\footnotesize,samepage=true,frame=single]
>>> dat_hyperlink['href']
'http://buildings.umwblogs.org/trinkle-hall/'
\end{Verbatim}

\subsection{The bottom-up approach}

Often, though, you quickly find yourself in a tangled mess as you try to dive
from the top of the document all the way to what you want. It's frequently
easier to search for the data you want using BeautifulSoup's search
facilities.

In the case above, suppose I wanted to create a DataFrame with two columns:
the roles Stephen plays, and the percentage of his total time each one takes.
The final form should look something like this:

\begin{Verbatim}[fontsize=\small,samepage=true,frame=none]
                role  perc
0  College professor  50.0
1                Dad  30.0
            ...
\end{Verbatim}


To find my way through the forest, I search through the HTML with developer
tools, locating one of the pieces of data I know is there. For instance, the
number \texttt{50}, which I can see on the page is one of the percentages.
When I click on this (or just search for it in a text file), I find myself in
the middle of the code in Figure~\ref{fig:bottomUpSearch}.

\begin{figure}[ht]
\centering
\begin{Verbatim}[fontsize=\footnotesize,samepage=true,frame=lines]
...
<table border=2 cellpadding=15>
<tr><td class="pic"><img src="a.jpg"></td>
    <td class="role">College professor</td>
    <td class="description">
    Stephen teaches computer science and data science to willing undergraduates,
    shaping them into the heroes of tomorrow.
    </td>
    <td>
    <ul>
    <li>Super rewarding</li>
    <li>Exciting</li>
    <li>Ever-changing</li>
    <li>Fun</li>
    </ul>
    </td>
    <td>
    <b>50%</b> of total time
    </td>
</tr>
...
\end{Verbatim}
\caption{Using a bottom-up approach -- locating the text ``\texttt{50}'' in
the HTML file and getting one's bearings regarding where it is in relation to
other elements.}
\label{fig:bottomUpSearch}
\end{figure}
\normalsize

Looks like the role name I want is the text under a \texttt{td} element whose
class is ``\texttt{role}''. And it looks like the percentage I want is
(*groan*) in a boldface (\texttt{b}) tag, underneath the last \texttt{td} tag
of the row, which has no \texttt{class}.

BeautifulSoup's \texttt{.find\_all()} method makes this slightly less of a
headache. To quickly grab all the \texttt{tr} elements in the document, I can
do this:

\begin{Verbatim}[fontsize=\footnotesize,samepage=true,frame=single]
trs = list(soup.find_all("tr"))
\end{Verbatim}

The variable \texttt{trs} is now a list of \textit{elements} (BeautifulSoup
``\texttt{Tag}'' objects) each of which I can examine for (say) \texttt{td}
children with a \texttt{role} \texttt{class} attribute. And you can call
\texttt{.find\_all()} on \textit{any} tag (to search from that point in the
hierarchy down), not just on the top-level \texttt{soup} element. So:

\begin{Verbatim}[fontsize=\footnotesize,samepage=true,frame=single]
roles = []
trs = soup.find_all("tr")
for tr in trs:
    tds = tr.find_all("td")
    for td in tds:
        if 'class' in td.attrs and 'role' in td['class']:
            ...
\end{Verbatim}

(A slight weirdness in that \texttt{if} statement: notice that I'm checking
whether \texttt{'role'} is \textbf{\texttt{in}} \texttt{td['class']} instead
of simply being \texttt{==} to \texttt{td['class']}. That's because in HTML,
the \texttt{class} attribute of an tag can have \textit{multiple}
(space-separated) values, not just one, and so \texttt{td['class']} is
actually a \textit{list}, not a single string.)

Now when we reach the body of that \texttt{if} statement, we have successfully
found one of the desired \texttt{td} attributes. How do we then get the text
under it? Just call \texttt{.get\_text()} on it, and then append it to the
list. This creates my desired list of roles that we can stick in my DataFrame.

\begin{Verbatim}[fontsize=\footnotesize,samepage=true,frame=single]
            ...
            roles.append(td.get_text())
\end{Verbatim}

To get the percentages, we just need a slightly different approach: first,
we're going to assume the percentage is in the \textit{last} \texttt{td} tag,
and second, we need to go inside that to retrieve the percentage from the
embedded boldface tag. Adding this to our \texttt{for} loop almost gets us
there:

\begin{Verbatim}[fontsize=\footnotesize,samepage=true,frame=single]
percs = []
...
    the_bold = tds[-1].find("b")
    percs.append(the_bold.get_text())
\end{Verbatim}

although we'd like to also strip that percentage sign off and convert it to an
actual numeric value, so we'll do this instead:

\begin{Verbatim}[fontsize=\footnotesize,samepage=true,frame=single]
percs = []
...
    the_bold = tds[-1].find("b")
    percs.append(float(the_bold.get_text()[:-1]))
\end{Verbatim}

and now we can create our DataFrame with both columns. (Note that
``\texttt{.find()}'' does the same thing as ``\texttt{.find\_all()}'' except
that it only grabs the first matching element rather than all of them.) The
complete code is shown in Figure~\ref{fig:getRolesComplete}.

\begin{figure}[ht]
\begin{Verbatim}[fontsize=\footnotesize,samepage=true,frame=single]
roles = []
percs = []

trs = soup.find_all("tr")
for tr in trs:
    tds = tr.find_all("td")
    for td in tds:
        if 'class' in td.attrs and 'role' in td['class']:
            roles.append(td.get_text())
    the_bold = tds[-1].find("b")
    percs.append(float(the_bold.get_text()[:-1]))

df = pd.DataFrame({'role':roles, 'perc':percs})[['role','perc']]
\end{Verbatim}
\caption{The complete code example to screen scrape the roles and percentages
from the web page in Figure~\ref{fig:allAboutStephen}
(p.~\pageref{fig:allAboutStephen}).}
\label{fig:getRolesComplete}
\end{figure}

\subsection{Other bs4 features}

The BeautifulSoup library has a myriad of alternatives for navigating and
searching the tree. Here are a few.

\subsubsection{Searching by id or class}

\texttt{.find()} and \texttt{.find\_all()} each have optional \texttt{id} and
\texttt{class\_} (note the trailing underscore!) parameters that can be used
to restrict a search to only elements with matching attributes.

\begin{Verbatim}[fontsize=\footnotesize,samepage=true,frame=single]
soup.find_all("td", class_="pic")
\end{Verbatim}

\subsubsection{Using the Python ``dot'' syntax}

If we know the names of the tags, we can navigate a structure by using a
Pandas-like ``dot'' syntax. Check it out:

\begin{Verbatim}[fontsize=\footnotesize,samepage=true,frame=single]
>>> soup.head.title
    <title>All About Stephen</title>
>>> soup.head.name
    title
>>> soup.head.title.string
    'All About Stephen'
>>> soup.body.table.td.img['src']
    'a.jpg'
\end{Verbatim}

\subsubsection{Accessing all attributes}

In addition to treating an element (``\texttt{Tag}'') as a dictionary to
retrieve its element values, you can also use the ``\texttt{.attrs}'' syntax to
get the dictionary directly:

\begin{Verbatim}[fontsize=\footnotesize,samepage=true,frame=single]
>>> soup.body.find("a").attrs
    {'href': 'http://buildings.umwblogs.org/trinkle-hall/'}
>>> soup.body.find("td").attrs
    {'class': ['pic']}
\end{Verbatim}

(As noted before, see how the value of the ``\texttt{class}'' attribute is a
list here.)

\subsubsection{Getting just text with \texttt{.get\_text()}}

The aforementioned \texttt{.get\_text()} method is actually a bit more
powerful than just extracting the contents of a single plain text element
underneath an element. It can also be used to get \textit{all the text} from
that point down in the hierarchy, ignoring all tags. Check it out:

\begin{Verbatim}[fontsize=\footnotesize,samepage=true,frame=single]
>>> soup.body.find("p").get_text()
    'Deep in the bowels of Trinkle Hall resides a\ncarbon-based life form
known to outsiders only as Stephen. This\nmysterious being fulfills several
roles not widely known to\nhumankind.'
\end{Verbatim}

\subsubsection{Finding more than one tag name at once}

You can pass a \textit{list} to \texttt{.find\_all()} to match more than one
type of tag in a search:

\begin{Verbatim}[fontsize=\footnotesize,samepage=true,frame=single]
soup.find_all(["td", "th", "p"])
\end{Verbatim}

And if you know how to use \textbf{regular expressions} (outside the scope of
this class) you can use them as arguments to \texttt{.find\_all()} and
friends:

\begin{Verbatim}[fontsize=\footnotesize,samepage=true,frame=single]
soup.find_all(re.compile("^q"))
soup.find_all(re.compile("r"))
soup.find_all(re.compile("s$"))
\end{Verbatim}

(finds all tags whose name begin with q, contain r, or end with s,
respectively.)

\subsubsection{CSS selectors}

Finally, if you're familiar with CSS selector syntax, you can do all that jazz
by calling the \texttt{.select()} method:

\begin{Verbatim}[fontsize=\footnotesize,samepage=true,frame=single]
soup.select("div#master, tr td.data")
\end{Verbatim}


